\documentclass[12pt,a4paper]{article}
\usepackage[T2A]{fontenc}
\usepackage[utf8]{inputenc}
\usepackage[russian]{babel}
\usepackage{amsmath,amssymb}
\usepackage[margin=2cm]{geometry}

\title{Решение системы методом Гаусса}
\date{}

\begin{document}
\maketitle


Даны матрицы A и B.
\[
\left(
\begin{array}{ccc}
1 & 1 & 2\\
2 & 3 & 1\\
3 & 2 & 4
\end{array}
\right),
\qquad
\left(
\begin{array}{c}
7\\
8\\
14
\end{array}
\right).
\]

Рассмотрим расширенную матрицу системы (A|B):
\[
\left(
\begin{array}{ccc|c}
1 & 1 & 2 & 7\\
2 & 3 & 1 & 8\\
3 & 2 & 4 & 14
\end{array}
\right).
\]

Приведём матрицу к ступенчатому виду методом Гаусса:
\[
\left(
\begin{array}{ccc|c}
1 & 1 & 2 & 7\\
2 & 3 & 1 & 8\\
3 & 2 & 4 & 14
\end{array}
\right)
\xrightarrow{\substack{R_2\leftarrow R_2-2R_1}}
\left(
\begin{array}{ccc|c}
1 & 1 & 2 & 7\\
0 & 1 & -3 & -6\\
3 & 2 & 4 & 14
\end{array}
\right)
\xrightarrow{\substack{R_3\leftarrow R_3-3R_1}}
\left(
\begin{array}{ccc|c}
1 & 1 & 2 & 7\\
0 & 1 & -3 & -6\\
0 & -1 & -2 & -7
\end{array}
\right)
\]

\[
\xrightarrow{R_3\leftarrow R_3+R_2}
\left(
\begin{array}{ccc|c}
1 & 1 & 2 & 7\\
0 & 1 & -3 & -6\\
0 & 0 & -5 & -13
\end{array}
\right)
\xrightarrow{R_3\leftarrow -R_3}
\left(
\begin{array}{ccc|c}
1 & 1 & 2 & 7\\
0 & 1 & -3 & -6\\
0 & 0 & 5 & 13
\end{array}
\right).
\]

В последней матрице все три строки ненулевые. Поэтому
\[
\operatorname{rank}(A)=\operatorname{rank}(A|B)=3,
\]

Так как
\[
\operatorname{rank}(A)=\operatorname{rank}(A|B)=3!=n, где n = 4,
\]
то система совместна и имеет бесконечное число решений.

\end{document}
то система совместна и имеет бесконечное число решений.

\end{document}
